\documentclass[11pt]{article}
\usepackage{amsmath,booktabs}
\usepackage{tikz}
\usepackage{fontspec}
\usepackage{unicode-math}
\usepackage{polyglossia}
\setmainlanguage{english}
\setotherlanguages{arabic,hindi}
\setmainfont{Latin Modern Roman}
\newfontfamily\arabicfont[Script=Arabic]{Amiri}
\newfontfamily\hindifont[Script=Devanagari]{NewComputerModern 10 Devanagari}
\newfontfamily\zhfont{FandolSong}
\newfontfamily\jafont{IPAexMincho}
\begin{document}
\section{Introduction}\label{sec:intro}
Every writer knows the rhythm.

\begin{figure}[h!]
\centering
\newcommand\petals{5}
\begin{tikzpicture}[scale=.6]
  \foreach \i in {1,...,\petals}
    \draw[fill=red!15, rotate=\i*360/\petals] (0,0) .. controls (.7,.45) and (1.7,.35) .. (2,0) .. controls (1.7,-.35) and (.7,-.45) .. (0,0);
\end{tikzpicture}
\caption{A figure.}\label{fig:live}
\end{figure}

\begin{table}[h!]
\centering
\begin{tabular}{lr}\toprule Language & Script\\\midrule Arabic & RTL\\ Hindi & Devanagari\\\bottomrule\end{tabular}
\caption{Scripts.}
\end{table}

\textlang{arabic}{مرحبا بالعالم}

\textlang{hindi}{नमस्ते दुनिया। क्षत्रिय, ज्ञान, श्रद्धा, कि की कु कू के कै।}

{\zhfont 你好，世界！实时排版。}

{\jafont こんにちは、世界。リアルタイム組版。}

Section~\ref{sec:method} describes the method~\cite{knuth84}, with $\sum_{i=1}^n i = \frac{n(n+1)}{2}$.

\newpage
\section{Method}\label{sec:method}
A document is a program.
\[ x = 1 \]

\begin{thebibliography}{9}
\bibitem{knuth84} D.~E. Knuth. \emph{The \TeX book}. 1984.
\end{thebibliography}
\end{document}
